\documentclass[11pt,a4paper]{article}

\usepackage[utf8]{inputenc}

\usepackage[T1]{fontenc}

\usepackage[polish]{babel}

\usepackage{lmodern}

\usepackage{geometry}

\geometry{margin=27mm}

\usepackage[hidelinks]{hyperref}

\title{Informacyjne braki danych i nieregularne interwały pomiaru przywiązania}

\author{Mateusz Pomianek}

\date{Szkic naukowy, wersja 1.0.0 | 28 listopada 2026}

\begin{document}

\maketitle

\begin{abstract}

Proponowana miara główna to obciążenie estymatora trendu względem wartości z generatora.

\end{abstract}

\section{Wprowadzenie}

Pomiary chwilowe ujawniają dynamikę doświadczenia w kontekście (Shiffman et al., 2008). Modele mieszanin wzorców pozwalają jawnie uzależnić wnioski od założeń o brakujących obserwacjach (Little, 1993).

\section{Cel i pytania badawcze}

Hipotezy robocze wymagają zatwierdzenia przed rozpoczęciem właściwej weryfikacji.

\paragraph{H1.} Jeżeli prawdopodobieństwo odpowiedzi maleje wraz z nieobserwowaną intensywnością reakcji i ta zależność wzmacnia się w czasie, analiza dostępnych pomiarów zaniży dodatni trend w porównaniu z prawdą generatora.

\paragraph{H2.} Przy brakach MAR i nieregularnych odstępach model z rzeczywistym czasem zmniejszy błąd estymacji trendu względem modelu indeksującego czas numerem pomiaru.

\section{Materiał i metody}

Pierwszorzędowy wskaźnik: obciążenie estymatora trendu względem wartości z generatora.

Porównanie symulacyjne wykorzysta obciążenie estymatora nachylenia i pokrycie 95-procentowych przedziałów. Analiza MNAR zmienia założoną wielkość różnicy między obserwowanymi i brakującymi wynikami oraz raportuje, kiedy odwraca się znak wniosku.

Przed badaniem z udziałem ludzi należy zatwierdzić wymagania etyczne, zasady zgody i ochronę danych.

\section{Plan analizy}

Zamrożenie estymandu, reguł oceny jakości i obsługi braków nastąpi przed dostępem do danych wynikowych. Raport obejmie oszacowanie kontrastu głównego z przedziałem niepewności i pełne mianowniki. Jeśli przewidziano hipotezy drugorzędowe, przed analizą należy ustalić sposób kontroli wielokrotności.

\section{Wyniki}

W niniejszym pliku nie umieszczono wyników empirycznych, statystyk ani nieprzeprowadzonych analiz. Sekcja wymaga uzupełnienia po uzyskaniu i weryfikacji materiału dowodowego.

\section{Dyskusja i ograniczenia}

Ograniczenia obejmują możliwe błędy pomiaru, dobór próby, alternatywne mechanizmy i ograniczoną możliwość uogólniania. Korelacja nie dowodzi przyczynowości; poprawność techniczna nie jest dowodem znaczenia klinicznego.

\section{Wnioski}

Hipotezy i sposób pomiaru stanowią podstawę oceny wykonalności oraz przygotowania właściwego protokołu.

\section*{Bibliografia}

{\setlength{\parindent}{0pt}
\setlength{\parskip}{5pt}
\hangindent=1.27cm
Little, R. J. A. (1993). Pattern-mixture models for multivariate incomplete data. \textit{Journal of the American Statistical Association, 88}(421), 125--134. \url{https://doi.org/10.1080/01621459.1993.10594302}\par
\hangindent=1.27cm
Shiffman, S., Stone, A. A., \& Hufford, M. R. (2008). Ecological momentary assessment. \textit{Annual Review of Clinical Psychology, 4}, 1--32. \url{https://doi.org/10.1146/annurev.clinpsy.3.022806.091415}\par
}

\end{document}
